\section{消融与分析}

\subsection{物理重排序的作用}

第一组消融移除了生成后的物理重排序，直接依据生成模型或代理模型输出选择候选结构。结果表明，这会显著增大名义性能与验证性能之间的差距，即便被选中的结构在视觉上看起来相当规整。这说明物理重排序并不是锦上添花的后处理，而是保证生成式逆向设计与真实电磁响应对齐的关键组成部分。

这一现象也解释了本文为何强调“物理一致”。生成模型本身并不能自动保证一致性，尤其当逆问题具有明显多解性、而训练样本又只能稀疏覆盖响应流形时更是如此。因此，闭环物理验证是将“生成式候选提出器”转化为“可靠逆向设计系统”的必要条件。

\subsection{统一响应条件的作用}

第二组消融将完整角度--频率--偏振目标替换为更简化的条件输入，例如单切片目标或压缩后的标量描述符。这些变体通常能够恢复部分算子行为，但在验证响应中难以同时保持角向结构、光谱稳定性和偏振纯度。角向畸变、光谱漂移和偏振泄漏在目标被过度压缩时都会更明显。

这一结果直接支持了本文的响应空间论点。若目标描述被压缩得过于简单，逆向模型就无法获得区分物理上不同可行解所需的信息。完整响应场条件的价值因此不仅体现在数值精度上，更体现在它维持了目标表述与器件物理之间的一致性。

\subsection{生成式与确定性逆向映射的比较}

最后，我们比较了确定性回归器与条件生成模型。确定性基线往往能给出平均意义上尚可的响应，但其解的多样性较弱，经过验证后的最佳性能通常也不如生成式模型。相比之下，生成模型能够针对同一目标输出一个候选解集，再由物理重排序挑选出验证保真度更高的结构。

对于自由形态傅里叶算子超表面而言，这一点尤为关键，因为逆问题的非唯一性是问题本身的属性，而不是训练中的噪声。生成式逆向设计的价值因此不仅在于“采样更多”，更在于它尊重了统一响应空间中一对多映射的本质。
